% Build twice with pdfLaTeX from this directory or the repository root.
\makeatletter
\def\input@path{{./}{../../../preamble/}{book2026-2027/preamble/}}
\makeatother
\input{preamble.tex}
\input{graphs/highlights.tex}
% \GTenHighlightsfalse % Uncomment to disable instructional highlights.
\input{graphs/specific_highlights.tex}
\input{graphs/G10_T2_bar_graphs.tex}

\geometry{letterpaper,margin=0.75in}
\hypersetup{pdftitle={Grade 10 Global Economics -- Term 2 Review Outline},hidelinks}
\linespread{1.1}
\setlength{\parskip}{4pt minus 2pt}
\setlength{\emergencystretch}{2em}
\setlist[enumerate]{leftmargin=*,topsep=3pt,itemsep=1pt,parsep=0pt}

\newcolumntype{L}{>{\raggedright\arraybackslash}X}
\newcolumntype{R}[1]{>{\raggedright\arraybackslash}p{#1}}
\newcommand{\criterionlink}[1]{\hyperref[crit:#1]{\textbf{\underline{#1}}}}
\newcommand{\backtomap}{\hfill\hyperref[sec:map]{\small\underline{Criterion map}}\par}
\newcommand{\GTenModelClue}[1]{\ifGTenHighlights\textbf{#1}\else#1\fi}

\begin{document}

\begin{center}
{\LARGE\bfseries Grade 10 Global Economics --- Term 2 Review Outline\par}
\smallskip
{\large 2026--2027\par}
\end{center}

\setcounter{tocdepth}{1}
\phantomsection
\label{toc}
\tableofcontents

\section{Introduction}

In Lesson 1 (C1--C4), you assign numbers to complete outcomes, group their probabilities, and begin using models for single events and fixed trial counts. In Lesson 2 (C5--C7), the observation rule tells you whether to count trials to a target or count events over a fixed exposure. In Lesson 3 (C8--C9), you interpret expected return and financial risk and decide whether a model fits the evidence.

\textbf{Use the same routine throughout the term.}
\begin{enumerate}
\item Define the variable in words and give its units. For a binary trial, name the event called success.
\item Identify what is fixed and what is random.
\item State the support (values with positive probability). Check the model assumptions or validate the supplied probability table.
\item Translate the wording into an event. Mark which endpoints are included.
\item Show the formula and substitution, check the result, and interpret it in context.
\end{enumerate}
Probabilities are unitless. Keep exact expressions or full precision in your working. Round final probabilities to four decimal places unless that would hide a small positive probability.

\par\nobreak
\secInicio

\clearpage
\section{Summarized Outline}\label{sec:map}
Select a criterion to jump to its worked example. The next page gives the event and model tables.

{\small\noindent
\begin{tabularx}{\linewidth}{@{}l c L L L@{}}
\toprule
\textbf{Criterion} & \textbf{Lesson} & \textbf{When asked to\ldots} & \textbf{Method} & \textbf{Show in your working}\\
\midrule
\criterionlink{C1} & 1 & Turn outcomes into a distribution & Map outcomes, then group by value & Outcome masses, assigned values, support, total mass\\
\criterionlink{C2} & 1 & Read or check a PMF or CDF & Accumulate masses or recover jumps & Included values, endpoint heights, validity checks\\
\criterionlink{C3} & 1 & Describe one yes/no event & Build a Bernoulli indicator & Meaning of 0 and 1, both masses, money mapping\\
\criterionlink{C4} & 1 & Count successes in fixed trials & Use binomial point masses, a sum, or a complement & Four conditions, \(n,p\), and arrangement count\\
\criterionlink{C5} & 2 & Wait for the first success & Use geometric point or tail probabilities & Failure--success sequence, cutoff, any conditioning\\
\criterionlink{C6} & 2 & Wait for success \(r\) & Fix the final success and count earlier arrangements & \(r,p\), final-trial condition, deadline sum or check\\
\criterionlink{C7} & 2 & Count events over an exposure & Convert units and use Poisson & Exposure, assumptions, \(\lambda\), included counts\\
\criterionlink{C8} & 3 & Compare expected gains and risk & Compute weighted moments and check downside & Calculation table, units, loss chance and worst value\\
\criterionlink{C9} & 3 & Choose or justify a model & Follow the observation rule and test assumptions & Support, parameters, and why an alternative fails\\
\bottomrule
\end{tabularx}\par
}

\newpage
\backtomap
\PracticeHighlightLegend{G10-T2-C1-C2-C3-C4-C5-C6-C7-C8-C9}
\subsection*{Quick reference for translating events}

Write \(p_X(a)=P(X=a)\) for the probability mass function (PMF) and
\(F_X(a)=P(X\leq a)\) for the cumulative distribution function (CDF).
The left limit \(F_X(a^-)=P(X<a)\) excludes any mass at \(a\).
Use \(T\) for a stopping trial and \(m\) for a deadline.

\begin{tabularx}{\linewidth}{@{}l l l L@{}}
\toprule
\textbf{Words} & \textbf{Event} & \textbf{Calculation} & \textbf{Included values if \(a=2\)}\\
\midrule
Exactly \(a\) & \(X=a\) & \(p_X(a)=F_X(a)-F_X(a^-)\) & \(\GTenHLValue{2}\)\\
Less than \(a\) & \(X<a\) & \(F_X(a^-)\) & \(\GTenHLValue{0}\)\\
At most \(a\) & \(X\leq a\) & \(F_X(a)\) & \(\GTenHLValue{0,2}\)\\
At least \(a\) & \(X\geq a\) & \(1-F_X(a^-)\) & \(\GTenHLValue{2,5}\)\\
More than \(a\) & \(X>a\) & \(1-F_X(a)\) & \(\GTenHLValue{5}\)\\
By trial \(m\) & \(T\leq m\) & \(F_T(m)\) & Includes trial \(m\)\\
After trial \(m\) & \(T>m\) & \(1-F_T(m)\) & Starts at trial \(m+1\)\\
\bottomrule
\end{tabularx}

The right column uses the \criterionlink{C2} score support \(\GTenHLValue{\{0,2,5\}}\) for the first five rows. With masses \(\GTenHLProbability{0.2,0.5,0.3}\), less than 2 has probability \(\GTenHLProbability{0.2000}\), while at most 2 has probability \(\GTenHLProbability{0.7000}\). A cutoff can fall in a gap. For example, \(1<X<5\) contains just \GTenHLValue{2}, so its probability is \(\GTenHLProbability{0.5000}\).

\subsection*{Quick reference for choosing a model}
Here \(\GTenHLProbability{p}\) is the success probability per trial and \(\GTenHLProbability{q}=1-p\). The listed supports assume \(0<p<1\) and \(\lambda>0\). Mass outside the support is zero. Both waiting-time models include the final successful trial.

{\small
\setlength{\tabcolsep}{4pt}
\noindent
\begin{tabularx}{\linewidth}{@{}R{0.14\linewidth}R{0.14\linewidth}R{0.14\linewidth}R{0.13\linewidth}L l@{}}
\toprule
\textbf{Model} & \textbf{Fixed} & \textbf{Variable} & \textbf{Support} & \textbf{Parameters and assumptions} & \textbf{Point mass}\\
\midrule
Bernoulli & \GTenModelClue{One observation} & Indicator \(\GTenHLValue{X}\) & \(\{0,1\}\) & \(p\). Two exclusive, exhaustive outcomes. & \(p^x q^{1-x}\)\\
\addlinespace
Binomial & \GTenModelClue{\(n\) trials} & Success count \(\GTenHLValue{X}\) & \(\{0,\ldots,n\}\) & \(n,p\). Independent binary trials, constant \(p\). & \(\binom nk p^k q^{n-k}\)\\
\addlinespace
Geometric & \GTenModelClue{First success} as target & Total trials \(\GTenHLValue{T}\) & \(\{1,2,\ldots\}\) & \(p\). Independent binary trials, constant \(p\). & \(q^{t-1}p\)\\
\addlinespace
Negative binomial & Target of \GTenModelClue{\(r\) successes} & Total trials \(\GTenHLValue{T}\) & \(\{r,r+1,\ldots\}\) & \(r,p\). Independent binary trials, constant \(p\). & \(\binom{t-1}{r-1}p^r q^{t-r}\)\\
\addlinespace
Poisson & \GTenModelClue{Exposure} (e.g., time or length) & Event count \(\GTenHLValue{X}\) & \(\{0,1,\ldots\}\) & \(\lambda\). Constant rate, independent disjoint counts, isolated events. & \(e^{-\lambda}\lambda^k/k!\)\\
\bottomrule
\end{tabularx}\par
}

Here \(\lambda=\text{rate}\times\text{exposure}\) is the mean count in the chosen exposure. The binomial coefficient \(\binom nk=n!/[k!(n-k)!]\) counts ways to choose \(k\) positions from \(n\). The factorial \(k!\) is the product of the positive integers through \(k\), with \(0!=1\).
\par\nobreak
\secInicio

\newpage
\section{Step-by-Step Outline}
\subsection{C1 --- From outcomes to a probability mass function}\label{crit:C1}
\textbf{Use this when} several complete outcomes can produce the same numerical result, or when a probability table has an unknown mass.

\begin{enumerate}
\item List the complete outcomes \(\GTenHLOutcome{\omega}\), assign each a value \(\GTenHLValue{X}(\GTenHLOutcome{\omega})\), and record its probability. Multiply across independent decisions when the context supplies independence.
\item List the supported values and their units. For each value \(x\), add the probabilities of all outcomes that map to it.
\item Set the sum of the masses equal to 1 to check or solve for unknowns. Check non-negativity and include any tail beyond the displayed values.
\end{enumerate}

\textbf{Notebook buyers.} Three visitors independently buy with probability \(\GTenHLProbability{0.4}\) each. \GTenHLOutcome{B} means buy and \GTenHLOutcome{N} means no purchase. In the table, each label records the decisions of all three visitors in order. The markers identify \GTenHLOutcome{outcome labels}, \GTenHLValue{assigned values and units}, and \GTenHLProbability{unitless probabilities}.

{\small\noindent
\begin{tabularx}{\linewidth}{@{}L c c c@{}}
\toprule
\GTenHLOutcome{Complete outcomes \(\omega\)} &
\(\GTenHLValue{x}=\GTenHLValue{X}(\GTenHLOutcome{\omega})\) \GTenHLValue{(buyers)} &
\GTenHLProbability{Mass of each outcome} &
\GTenHLProbability{Grouped mass \(p_X(x)\)}\\
\midrule
\GTenHLOutcome{NNN} & \GTenHLValue{0} & \(0.6^3=\GTenHLProbability{0.216}\) & \(\GTenHLProbability{0.2160}\)\\
\GTenHLOutcome{BNN, NBN, NNB} & \GTenHLValue{1} & \(0.4(0.6)^2=\GTenHLProbability{0.144}\) & \(3(0.144)=\GTenHLProbability{0.4320}\)\\
\GTenHLOutcome{BBN, BNB, NBB} & \GTenHLValue{2} & \((0.4)^2(0.6)=\GTenHLProbability{0.096}\) & \(3(0.096)=\GTenHLProbability{0.2880}\)\\
\GTenHLOutcome{BBB} & \GTenHLValue{3} & \(0.4^3=\GTenHLProbability{0.064}\) & \(\GTenHLProbability{0.0640}\)\\
\bottomrule
\end{tabularx}\par
}

The support is \GTenHLValue{\(\{0,1,2,3\}\) buyers}. For instance, \GTenHLOutcome{BNN} and \GTenHLOutcome{NBN} describe different buyers but both contribute to \(p_X(\GTenHLValue{1})\). Exactly one buyer has probability \(\GTenHLProbability{0.4320}\). If the last mass were missing, \(\GTenHLParameter{p_X(3)}=1-(0.216+0.432+0.288)=\GTenHLParameter{0.0640}\), the chance all three buy.

\textbf{Normalizing weights.} If refunds of \GTenHLValue{\(0,2,5\) USD} have masses \(\GTenHLParameter{c,2c,3c}\), then \(6c=1\), so \(\GTenHLParameter{c}=\GTenHLParameter{1/6}\). The resulting masses \(\GTenHLProbability{1/6,1/3,1/2}\) are positive and total 1. A refund of at least 2 USD has probability \(1/3+1/2=5/6\approx\GTenHLProbability{0.8333}\).

\textbf{Infinite support.} If each independent call sells with probability \(\GTenHLProbability{1/2}\), the first sale on call \(t\) has mass \(p_T(t)=\GTenHLProbability{2^{-t}}\) for \(\GTenHLValue{t}=\GTenHLValue{1,2,\ldots}\). The first \(m\) terms sum to \(1-2^{-m}\), which tends to 1. Beyond call 5 the omitted mass is \(P(\GTenHLValue{T}>\GTenHLValue{5})=2^{-5}=\GTenHLProbability{1/32}=0.03125\).

\noindent
\begin{minipage}[t]{.48\linewidth}
\centering
\GTenPMFGraph[width=\linewidth,height=4.25cm,xlabel={\GTenHLValue{\(x\) (buyers)}},ylabel={\GTenHLProbability{Probability}},xhighlight=GTenHLValue,yhighlight=GTenHLProbability,color=GTenHLProbability!32]{0,1,2,3}{.216,.432,.288,.064}
\par\small \textbf{Buyer-count PMF.} Each height is the grouped probability in the last column above.
\end{minipage}\hfill
\begin{minipage}[t]{.48\linewidth}
\centering
\GTenPMFGraph[width=\linewidth,height=4.25cm,xlabel={\GTenHLValue{\(t\) (calls)}},ylabel={\GTenHLProbability{Probability}},xhighlight=GTenHLValue,yhighlight=GTenHLProbability,color=GTenHLProbability!32,tail=1/32]{1,2,3,4,5}{1/2,1/4,1/8,1/16,1/32}
\par\small \textbf{First-sale PMF, first five calls.} The bars total \(31/32\). Calls \(6,7,\ldots\) carry the remaining \(1/32\).
\end{minipage}
\par\nobreak\vspace{6pt}
\secInicio[][sec:map]

\newpage
\subsection{C2 --- Read, build, and verify a CDF}\label{crit:C2}
\textbf{Use this when} you are given individual masses or a cumulative table or graph, especially when the question changes an endpoint.

\begin{enumerate}
\item Verify that every PMF mass is non-negative and the total is 1.
\item Build \(F_X(\GTenHLValue{x})=\sum_{u\leq x}\GTenHLProbability{p_X(u)}\) by accumulating the masses in order. Keep the height constant between supported values.
\item To go back to the PMF, subtract the height just before a jump from the height at the jump. Check that the CDF never decreases, is right-continuous, and tends to 0 on the far left and 1 on the far right.
\item Mark the included values. Add their masses, subtract CDF heights, or use a complement.
\end{enumerate}

\textbf{Scores on irregular support.} \(\GTenHLValue{X}\) has support \(\GTenHLValue{\{0,2,5\}}\) in \GTenHLValue{score points}. Its masses are positive and \(0.2+0.5+0.3=1\). For any real cutoff \(\GTenHLValue{x}\),
\[
\begin{array}{l|ccc}
\GTenHLValue{x} & \GTenHLValue{0} & \GTenHLValue{2} & \GTenHLValue{5}\\ \hline
\GTenHLProbability{p_X(x)} & \GTenHLProbability{0.2} & \GTenHLProbability{0.5} & \GTenHLProbability{0.3}\\
\GTenHLProbability{F_X(x)} & \GTenHLProbability{0.2} & \GTenHLProbability{0.7} & \GTenHLProbability{1}
\end{array}
\qquad
F_X(x)=
\begin{cases}
0 & x<0,\\
\GTenHLProbability{0.2} & 0\leq x<2,\\
0.2+0.5=\GTenHLProbability{0.7} & 2\leq x<5,\\
0.7+0.3=\GTenHLProbability{1} & x\geq5.
\end{cases}
\]
Reading the jumps back gives \(p_X(0)=0.2-0=0.2000\), \(p_X(2)=0.7-0.2=0.5000\), and \(p_X(5)=1-0.7=0.3000\).

\pgfplotsset{score plot/.style={
  width=.44\linewidth,height=4.7cm,xmin=-1,xmax=6,ymin=-.07,ymax=1.1,
  xtick={0,2,5},xlabel={\GTenHLValue{\(x\) (score points)}},tick label style={font=\small},
  title style={font=\small},axis lines=box,ymajorgrids,grid style={gray!25}
}}
\ifGTenHighlights
\pgfplotsset{score plot/.append style={
  xticklabel style={gten marker=GTenHLValue,inner sep=1pt},
  yticklabel style={gten marker=GTenHLProbability,inner sep=1pt}
}}
\fi
\begin{center}
\begin{tikzpicture}
\begin{axis}[score plot,ytick={0,.2,.3,.5,.7,1},ylabel={\(\GTenHLProbability{p_X(x)}\)},title={PMF --- individual masses}]
\addplot[black,thick,ycomb,mark=*] coordinates {(0,.2) (2,.5) (5,.3)};
\end{axis}
\end{tikzpicture}\hfill
\begin{tikzpicture}
\begin{axis}[score plot,ytick={0,.2,.7,1},ylabel={\(\GTenHLProbability{F_X(x)}\)},title={CDF --- accumulated mass}]
\addplot[black,thick] coordinates {(-1,0) (0,0)};
\addplot[black,thick] coordinates {(0,.2) (2,.2)};
\addplot[black,thick] coordinates {(2,.7) (5,.7)};
\addplot[black,thick] coordinates {(5,1) (6,1)};
\addplot[black,only marks,mark=*] coordinates {(0,.2) (2,.7) (5,1)};
\addplot[black,only marks,mark=*,mark options={fill=white}] coordinates {(0,0) (2,.2) (5,.7)};
\end{axis}
\end{tikzpicture}
\end{center}
\textit{The filled dot gives the CDF value at each jump. The open dot gives its left limit. The two outer steps continue horizontally.}

\[
\begin{aligned}
P(0<X\leq5)&=F_X(5)-F_X(0)=1-0.2=\GTenHLProbability{0.8000},\\
P(X<\GTenHLValue{5})&=F_X(5^-)=\GTenHLProbability{0.7000},\qquad P(X\leq\GTenHLValue{5})=F_X(5)=\GTenHLProbability{1.0000},\\
P(X>2)&=1-F_X(2)=1-0.7=\GTenHLProbability{0.3000}.
\end{aligned}
\]
For \(P(a\leq X\leq b)\), subtract \(F_X(a^-)\) from \(F_X(b)\) to retain the mass at \(a\). Here a strict cutoff at 5 excludes the \(0.3000\) chance of scoring exactly 5.

\textbf{Reject invalid candidates.} Masses \(0.6,0.5,-0.1\) total 1 but contain a negative mass. Masses \(0.2,0.3,0.3\) total only \(0.8\). Reject a CDF that steps down, uses the open-dot height as its value at a jump, or levels off at \(0.9\) forever on the right.
\par\nobreak
\secInicio[][sec:map]

\newpage
\subsection{C3 --- Bernoulli indicators}\label{crit:C3}
\textbf{Use this when} \GTenModelClue{one observation} has \GTenModelClue{two mutually exclusive outcomes} that cover all possibilities, such as \GTenHLOutcome{default within one year} or \GTenHLOutcome{no default within one year}.

\begin{enumerate}
\item Define the event \(A\) and set \(\GTenHLValue{X}=\GTenHLValue{1}\) if \(A\) occurs and \(\GTenHLValue{X}=\GTenHLValue{0}\) otherwise. The indicator is unitless.
\item Show \(p_X(1)=p\) and \(p_X(0)=1-p\). To indicate the opposite event, set \(Z=1-X\) and swap the two masses.
\item Define a separate monetary variable if needed. Match each monetary outcome with its indicator value.
\end{enumerate}

\textbf{One loan default.} A loan has one-year default probability \(\GTenHLProbability{p}=\GTenHLProbability{0.04}\). Let \(X=\GTenHLValue{1}\) for default, so the support is \(\GTenHLValue{\{0,1\}}\). A default costs \GTenHLThird{800 USD}, giving loss \(\GTenHLThird{L}=800X\) in \GTenHLThird{USD}.
\[
\begin{array}{c|ccc|l}
\GTenHLOutcome{\text{Outcome}} & \GTenHLValue{X} & \GTenHLSecond{Z}=1-X & \GTenHLThird{L\text{ (USD)}} & \GTenHLProbability{\text{Probability}}\\ \hline
\GTenHLOutcome{\text{Default}} & \GTenHLValue{1} & \GTenHLSecond{0} & \GTenHLThird{800} & p=\GTenHLProbability{0.0400}\\
\GTenHLOutcome{\text{No default}} & \GTenHLValue{0} & \GTenHLSecond{1} & \GTenHLThird{0} & 1-0.04=\GTenHLProbability{0.9600}
\end{array}
\]
Thus \(P(\GTenHLSecond{Z}=\GTenHLSecond{1})=\GTenHLProbability{0.9600}\) and \(P(\GTenHLThird{L}=\GTenHLThird{800})=\GTenHLProbability{0.0400}\). The word \emph{success} names the event coded 1, which here triggers a loss.
\par\nobreak
\secInicio[][sec:map]

\subsection{C4 --- Binomial success counts}\label{crit:C4}
\textbf{Use this when} you \GTenModelClue{count successes} in a \GTenModelClue{fixed number of trials}.

\begin{enumerate}
\item Verify all four conditions by checking fixed \(n\), two outcomes per trial, independence, and constant success probability \(\GTenHLProbability{p}\). Define \(\GTenHLValue{X}\) as the success count with support \(\GTenHLValue{\{0,\ldots,n\}}\) for \(0<p<1\).
\item For exactly \(k\), multiply one arrangement's probability \(p^k(1-p)^{n-k}\) by \(\binom nk\), the number of ways to choose the \(k\) success positions.
\item For a range, add one mass per included count. For at least one, subtract the zero-success mass from 1.
\end{enumerate}

\textbf{Five independent loans.} Observe all five for one year. Each defaults with probability \(\GTenHLProbability{0.1}\), so \(\GTenHLValue{X}\sim\operatorname{Bin}(5,0.1)\) counts defaults. Exactly one default can occur in any of five positions.
\[
\begin{aligned}
P(X=\GTenHLValue{1})&=\binom51(0.1)(0.9)^4=0.32805\approx\GTenHLProbability{0.3281},\\
P(X\leq\GTenHLValue{1})&=P(X=0)+P(X=1)=(0.9)^5+5(0.1)(0.9)^4\approx\GTenHLProbability{0.9185},\\
P(X\geq\GTenHLValue{1})&=1-P(X=0)=1-(0.9)^5\approx\GTenHLProbability{0.4095}.
\end{aligned}
\]
About \(40.95\%\) of such portfolios would have one or more defaults under these assumptions. If a market shock makes the loans default together, a default on one loan changes what you know about the others. The independence calculation then needs a different joint model (see \criterionlink{C9}).
\par\nobreak
\secInicio[][sec:map]

\newpage
\subsection{C5 --- Geometric waiting to the first success}\label{crit:C5}
\textbf{Use this when} independent binary trials with constant success probability \(\GTenHLProbability{p}\) continue \GTenModelClue{until the first success}.

\begin{enumerate}
\item Define \(\GTenHLValue{T}\) as \GTenHLValue{total trials} through the first success, with support \(\GTenHLValue{\{1,2,\ldots\}}\) for \(0<p<1\). Set \(q=1-p\).
\item For \emph{on trial \(t\)}, write \(t-1\) failures followed by one success, giving \(\GTenHLProbability{q^{t-1}p}\).
\item For \emph{after trial \(m\)}, all first \(m\) trials fail, giving \(q^m\). For \emph{by trial \(m\)}, use \(1-q^m\). Convert money or failures to a total-trial cutoff first.
\end{enumerate}

\textbf{A window and a changed starting point.} Suppose independent calls sell with \(p=\GTenHLProbability{0.3}\). First sale on call 3 is \GTenHLOutcome{FFS}, so \(P(T=\GTenHLValue{3})=(0.7)^2(0.3)=\GTenHLProbability{0.1470}\). After call 3 has probability \(P(T>\GTenHLValue{3})=(0.7)^3=\GTenHLProbability{0.3430}\), and by call 3 has probability \(P(T\leq\GTenHLValue{3})=1-(0.7)^3=\GTenHLProbability{0.6570}\).
\[
\begin{aligned}
P(2<T\leq5)&=P(T>2)-P(T>5)=0.7^2-0.7^5\approx\GTenHLProbability{0.3219},\\
P(T\leq5\mid T>2)&=\frac{0.7^2-0.7^5}{0.7^2}=1-0.7^3=\GTenHLProbability{0.6570}.
\end{aligned}
\]
From the start, the window requires two initial failures. Given those failures, the next three calls retain their original chances. With \(\GTenHLSecond{F}=\GTenHLValue{T}-1\) failures before the sale, \(\GTenHLSecond{F}=\GTenHLSecond{2}\) is the event \(\GTenHLValue{T}=\GTenHLValue{3}\).

\textbf{Budget variant.} At 3 USD per call, a budget of 12 USD buys at most \(\lfloor12/3\rfloor=\GTenHLValue{4}\) calls. With the budget example's \(p=0.25\), the chance of making the first sale within budget is \(P(T\leq\GTenHLValue{4})=1-(0.75)^4\approx\GTenHLProbability{0.6836}\).
\par\nobreak\vspace{6pt}
\secInicio[][sec:map]

\subsection{C6 --- Negative binomial waiting to the target success}\label{crit:C6}
\textbf{Use this when} independent binary trials with constant \(\GTenHLProbability{p}\) continue \GTenModelClue{until a fixed target of \(r\) successes}.

\begin{enumerate}
\item Define \(\GTenHLValue{T}\) as \GTenHLValue{total trials} through success \(r\). Its support starts at \(r\). For \(T=\GTenHLValue{t}\), fix trial \(t\) as a success.
\item Choose \(r-1\) success positions among the first \(t-1\) trials. Multiply the resulting \(\binom{t-1}{r-1}\) arrangements by \(p^r(1-p)^{t-r}\).
\item For a deadline \(m\), sum the masses from \(t=r\) to \(m\). Check by counting at least \(r\) successes in a fixed window of \(m\) trials. If \(m<r\), the deadline probability is zero.
\end{enumerate}

\textbf{Third sale with \(p=0.4\).} \(\GTenHLValue{T}\) counts \GTenHLValue{calls}, with support \(\GTenHLValue{\{3,4,\ldots\}}\). The third sale on call 5 requires two sales among the first four calls and a sale on call 5. Thus \(P(T=\GTenHLValue{5})=\binom42(0.4)^3(0.6)^2=0.13824\approx\GTenHLProbability{0.1382}\).
\[
\begin{aligned}
P(\GTenHLValue{T}\leq\GTenHLValue{6})
 &= (0.4)^3[1+3(0.6)+6(0.6)^2+10(0.6)^3]
 =\GTenHLProbability{0.45568}\approx0.4557,\\
P(\GTenHLThird{B}\geq\GTenHLThird{3})
 &= 1-\sum_{j=0}^{2}\binom6j(0.4)^j(0.6)^{6-j}
 =1-0.54432=\GTenHLProbability{0.45568},\quad \GTenHLThird{B}\sim\operatorname{Bin}(6,0.4).
\end{aligned}
\]
Both calculations give the chance of reaching the target within six calls. In the binomial check, imagine \GTenModelClue{observing all six calls} even if the target is reached early. Failures satisfy \(\GTenHLSecond{F}=\GTenHLValue{T}-r\), so \(T=\GTenHLValue{5}\) gives \(F=\GTenHLSecond{2}\). With \(r=1\), the coefficient becomes \(\binom{t-1}{0}=1\) and the PMF becomes geometric.
\par\nobreak
\secInicio[][sec:map]

\newpage
\subsection{C7 --- Poisson counts over fixed exposure}\label{crit:C7}
\textbf{Use this when} you \GTenModelClue{count events} over a \GTenModelClue{fixed time, length, area, or other exposure}.

\begin{enumerate}
\item Define the count \(\GTenHLValue{X}\) and exposure. Check a \GTenModelClue{constant rate}, \GTenModelClue{independent counts in disjoint intervals}, and \GTenModelClue{isolated events} (occurring singly).
\item Convert the rate and exposure to compatible units. Multiply to get the mean count \(\GTenHLParameter{\lambda}\). The support is \(\GTenHLValue{\{0,1,2,\ldots\}}\) for \(\lambda>0\).
\item Use \(e^{-\lambda}\lambda^k/k!\) for exactly \(k\). Sum through \(k\) for at most \(k\), or subtract the sum through \(k-1\) from 1 for at least \(k\).
\end{enumerate}

\textbf{Arrivals over 150 minutes.} Assume isolated arrivals at a constant rate of \(0.8\) per hour and independent counts in disjoint intervals. Convert \(150\) minutes to \(150/60=2.5\) hours, so \(\GTenHLParameter{\lambda}=(0.8\text{ arrivals/hour})(2.5\text{ hours})=\GTenHLParameter{2}\) arrivals.
\[
\begin{aligned}
P(X=\GTenHLValue{2})&=e^{-2}\frac{2^2}{2!}=2e^{-2}\approx\GTenHLProbability{0.2707},\\
P(X\leq\GTenHLValue{2})&=e^{-2}\left(1+2+\frac{2^2}{2!}\right)=5e^{-2}\approx\GTenHLProbability{0.6767},\\
P(X\geq\GTenHLValue{3})&=1-5e^{-2}\approx\GTenHLProbability{0.3233}.
\end{aligned}
\]
About \(32.33\%\) of such intervals would contain at least three arrivals.

\textbf{Cost threshold with the same mean.} Suppose claims arrive under the Poisson assumptions at \(0.4\) claims per day and each costs 100 USD. Over five days the mean is again \(\GTenHLParameter{\lambda}=0.4(5)=\GTenHLParameter{2}\) claims. If \(\GTenHLSecond{L}=100\GTenHLValue{X}\) is total cost in \GTenHLSecond{USD}, then \(\GTenHLSecond{L}\geq\GTenHLSecond{300}\) means \(\GTenHLValue{X}\geq\GTenHLValue{3}\). The same complement gives \(P(L\geq300)=1-5e^{-2}\approx\GTenHLProbability{0.3233}\). Show the conversion from the money threshold to the count before using the PMF.

\textbf{When events cluster.} If a market shock triggers many defaults together, the isolated-event assumption fails for that default count. To calculate a tail probability, you need information about the frequency and size of the clusters. See \criterionlink{C9} for how a shared shock can change a portfolio probability.
\par\nobreak
\secInicio[][sec:map]

\newpage
\subsection{C8 --- Expected value, variance, and financial risk}\label{crit:C8}
\textbf{Use this when} a PMF gives uncertain gains, costs, or counts and you need to summarize the average and the spread.

\begin{enumerate}
\item Validate the PMF. For each supported value, show \(\GTenHLValue{x}\), \(\GTenHLProbability{p_X(x)}\), \(xp_X(x)\), and \(x^2p_X(x)\).
\item Sum the third column for the mean \(\GTenHLParameter{\mu}=\GTenHLParameter{E(X)}\) and the fourth for \(\GTenHLParameter{E(X^2)}\). Subtract \(\mu^2\) to get variance, then take its square root for standard deviation.
\item For \(\GTenHLSecond{Y}=a\GTenHLValue{X}+b\), use \(E(Y)=aE(X)+b\), \(\operatorname{Var}(Y)=a^2\operatorname{Var}(X)\), and \(\sigma_Y=|a|\sigma_X\). Compare loss probability and worst outcome as well as the spread.
\end{enumerate}

\textbf{Claims and a reserve balance.} \(\GTenHLValue{X}\) counts \GTenHLValue{daily claims}. A firm starts with 100 USD and pays 30 USD per claim, so \(\GTenHLSecond{Y}=100-30\GTenHLValue{X}\) is the \GTenHLSecond{remaining reserve in USD}.
\begin{center}
\begin{tabular}{@{}lrrr@{}}
\toprule
\GTenHLValue{\(x\) (claims)} & \(\GTenHLProbability{p_X(x)}\) & \(xp_X(x)\) (claims) & \(x^2p_X(x)\) (claims\(^2\))\\
\midrule
\GTenHLValue{0} & \GTenHLProbability{0.6} & 0 & 0\\
\GTenHLValue{1} & \GTenHLProbability{0.3} & \(1(0.3)=0.3\) & \(1^2(0.3)=0.3\)\\
\GTenHLValue{2} & \GTenHLProbability{0.1} & \(\GTenHLValue{2}(\GTenHLProbability{0.1})=0.2\) & \(\GTenHLValue{2}^2(\GTenHLProbability{0.1})=0.4\)\\
\midrule
Total & 1 & \GTenHLParameter{0.5} & \GTenHLParameter{0.7}\\
\bottomrule
\end{tabular}
\end{center}
Variance is the weighted average of squared deviations from the mean. Expanding the square and using \(\sum_x p_X(x)=1\) gives the shortcut:
\[
\sum_x (x-\mu)^2p_X(x)
= E(X^2)-2\mu E(X)+\mu^2
= E(X^2)-\mu^2.
\]
The mean is \(\GTenHLParameter{0.5}\) claims per day. The variance is \(0.7-0.5^2=\GTenHLParameter{0.45}\) claims\(^2\) and \(\sigma_X=\sqrt{0.45}\approx\GTenHLParameter{0.6708}\) claims. For the reserve,
\[
\begin{aligned}
E(Y)&=100-30(0.5)=\GTenHLParameter{85}\text{ USD},\\
\operatorname{Var}(Y)&=(-30)^2(0.45)=\GTenHLParameter{405}\text{ USD}^2,\qquad
\sigma_Y=\sqrt{405}\approx\GTenHLParameter{20.1246}\text{ USD}.
\end{aligned}
\]
Over many comparable days the average remaining reserve approaches 85 USD. On a single day the supported balances are \GTenHLSecond{100, 70, and 40 USD}.

\textbf{Equal variance can hide different losses.} Investment A gives \(\GTenHLValue{-10}\) or \GTenHLValue{\(10\) USD} with probability \(\GTenHLProbability{0.5}\) each. B gives \GTenHLSecond{\(-20,0,20\) USD} with masses \(\GTenHLProbability{0.125,0.75,0.125}\). The opposite gains cancel, giving both means 0 USD. Their second moments are \(2(10^2)(0.5)=100\) and \(2(20^2)(0.125)=100\) USD\(^2\), so both variances are \(100\) USD\(^2\).
\begin{center}
\begin{tabular}{@{}lcc@{}}
\toprule
 & \textbf{Investment A} & \textbf{Investment B}\\
\midrule
Probability of a loss & \(\GTenHLProbability{0.5000}\) & \(\GTenHLProbability{0.1250}\)\\
Worst gain (USD) & \(\GTenHLValue{-10}\) & \(\GTenHLSecond{-20}\)\\
\bottomrule
\end{tabular}
\end{center}
B loses less often, but its worst loss is twice as large.
\par\nobreak
\secInicio[][sec:map]

\newpage
\subsection{C9 --- Select and justify the probability model}\label{crit:C9}
\textbf{Use this when} several models sound plausible or the question asks whether enough information is available.

\begin{enumerate}
\item Define \GTenModelClue{what the variable measures} and \GTenModelClue{what the observation rule fixes}.
\item Follow the chart to a candidate. State its support, units, parameters, and assumptions. Explain which feature rules out a plausible competing model.
\item Calculate with the justified model. If a required assumption is unresolved, show which probability depends on the missing information.
\end{enumerate}

\begin{center}
\begin{tikzpicture}[
  question/.style={draw,rounded corners,text width=6.3cm,align=center,font=\small,minimum height=7mm},
  answer/.style={draw,text width=5.2cm,align=center,font=\small,minimum height=7mm},
  >=Stealth
]
\node[question] (q1) at (0,0) {\GTenModelClue{One binary indicator}?};
\node[answer] (a1) at (7.6,0) {Bernoulli \(\bigl(p\bigr)\)};
\node[question] (q2) at (0,-1.05) {Success count in \GTenModelClue{fixed \(n\) trials}?};
\node[answer] (a2) at (7.6,-1.05) {Binomial \(\bigl(n,p\bigr)\)};
\node[question] (q3) at (0,-2.10) {Total trials \GTenModelClue{until the first success}?};
\node[answer] (a3) at (7.6,-2.10) {Geometric \(\bigl(p\bigr)\)};
\node[question] (q4) at (0,-3.15) {Total trials \GTenModelClue{until success \(r>1\)}?};
\node[answer] (a4) at (7.6,-3.15) {Negative binomial \(\bigl(r,p\bigr)\)};
\node[question] (q5) at (0,-4.20) {Event count over \GTenModelClue{fixed exposure}?};
\node[answer] (a5) at (7.6,-4.20) {Poisson \(\bigl(\lambda\bigr)\)};
\node[question] (q6) at (0,-5.45) {Get more information or try another model};
\foreach \i in {1,2,3,4,5}{\draw[->] (q\i) -- node[above,font=\small]{yes} (a\i);}
\foreach \i/\j in {1/2,2/3,3/4,4/5,5/6}{\draw[->] (q\i) -- node[right,font=\small]{no} (q\j);}
\end{tikzpicture}

\small \textit{A yes branch gives a candidate. Verify its assumptions using the model table.}
\end{center}

\textbf{On-time invoices.} \GTenModelClue{Observe all five invoices}. Each is independently paid on time with probability \(\GTenHLProbability{0.8}\). \(\GTenHLValue{X}\) counts on-time payments, with support \GTenHLValue{\(\{0,\ldots,5\}\) payments}. A \GTenModelClue{fixed \(n=5\)}, \GTenModelClue{binary outcomes}, \GTenModelClue{independence}, and \GTenModelClue{constant} \(p=\GTenHLProbability{0.8}\) justify binomial. Geometric would measure the trials \GTenModelClue{to the first on-time payment}.
\[
P(X=\GTenHLValue{4})=\binom54(0.8)^4(0.2)=5(0.4096)(0.2)=\GTenHLProbability{0.4096}.
\]
This is the chance that exactly one invoice is late.

\textbf{Shared default shock.} Each of five loans has one-year default probability \(\GTenHLProbability{0.1}\), so one loan's indicator has a Bernoulli PMF. The count \(\GTenHLValue{D}\) is between 0 and 5 defaults. Its support and aggregate probabilities depend on \GTenModelClue{how the loans default together}.
\begin{itemize}
\item If defaults are \GTenModelClue{independent}, \(D\sim\operatorname{Bin}(5,0.1)\) and \(P(D=\GTenHLValue{5})=(0.1)^5=\GTenHLProbability{0.00001}\).
\item If a shock with probability \(\GTenHLProbability{0.1}\) causes \GTenModelClue{all five to default together} and otherwise none default, the support is \(\GTenHLValue{\{0,5\}}\) and \(P(D=\GTenHLValue{5})=\GTenHLProbability{0.1000}\).
\end{itemize}
Both cases give each loan the same default probability. To assign a chance to all five defaulting, you need a joint model for the portfolio.
\par\nobreak
\secInicio[][sec:map]

\end{document}
